\chapter{Technológiai háttér}

A WaccApp technológiai hátterének bemutatásakor a hangsúly a frontend működését meghatározó kliensoldali megoldásokon van. A rendszer nem keretrendszerre épülő egyoldalas alkalmazásként készült, hanem több HTML-nézetből, CSS-ből és JavaScriptből álló webes frontendként. Ez a megközelítés a projekt méretéhez és felépítéséhez illeszkedett, mert a fő funkciók eltérő felhasználói helyzeteket képviselnek. Más felületet igényel egy üzenet elküldése, egy beszélgetés megtekintése, egy kérdőív szerkesztése, egy médiafájl csatolása vagy egy szűrés lefuttatása.

A technológiai fejezet célja nem a teljes backend részletes bemutatása, hanem annak tisztázása, hogy a frontend milyen alapokra épül, és ezek az alapok hogyan teszik lehetővé a későbbi modulok működését. A WaccApp-ban a HTML adja az egyes nézetek szerkezetét, a CSS biztosítja az elrendezést, a vizuális egységet és a reszponzív megjelenést, a JavaScript pedig a felhasználói események, az adatgyűjtés, az API-hívások és a nézetfrissítések kezeléséért felel. A frontend így nem önmagában álló statikus oldalak gyűjteménye, hanem olyan kliensoldali réteg, amely a háttérrendszerrel együtt teszi használhatóvá a kommunikációs funkciókat.

Ebben a fejezetben ezért a többnézetes frontend technológiai alapjai, a REST API-alapú kliens–szerver kapcsolat, a kliensoldali adatkezelés, az aszinkron működés, a validáció és a futtatási szempontok kerülnek előtérbe. A részletes modulbemutatás későbbi fejezetekben szerepel, itt elsősorban az a kérdés, milyen technológiai döntések tették lehetővé a WaccApp frontendjének működését. [R9], [R10], [R12]

\section{A rendszer frontend technológiai alapjai}

A WaccApp frontendjének alapját a több különálló HTML-nézetből felépülő kliensoldali szerkezet adja. A rendszer nem klasszikus egyoldalas alkalmazásként működik, ahol minden nézetváltást egy központi frontend keretrendszer kezel, hanem funkciók szerint elkülönített oldalakból áll. Ez a döntés azért volt indokolt, mert a rendszer fő feladatai nem ugyanazt a felhasználói folyamatot ismétlik. Egy gyors üzenetküldés, egy beszélgetés visszanézése, egy kérdőív összeállítása vagy egy szűrési eredmény értelmezése eltérő felületi szervezést kíván.

A rendszerben az \path{index.html} a küldési műveletek kiindulópontja, a \path{conv.html} a beszélgetések követéséhez, a \path{filter.html} az adatok visszakereséséhez, a kérdőíves és sablonkezelő oldalak pedig a kérdések, válaszok és sablonos folyamatok kezeléséhez kapcsolódnak. Ezeknek a nézeteknek a részletes szerepét a frontendarchitektúrát bemutató fejezet tárgyalja, technológiai szempontból azonban már itt fontos kiemelni, hogy a többoldalas felépítés nem véletlenszerű fájlszerkezet, hanem a funkciók elkülönítését támogató döntés.

Ebben a felépítésben a HTML az oldalak szerkezeti vázát adja. Itt jelennek meg azok az űrlapok, gombok, tartalmi blokkok és listaelemek, amelyekre a felhasználó közvetlenül reagál. A CSS feladata az, hogy ezek az elemek átlátható elrendezést, következetes vizuális hierarchiát és különböző kijelzőméreteken is használható megjelenést kapjanak. A JavaScript ehhez képest a működési réteget biztosítja: figyeli a felhasználói eseményeket, összegyűjti az adatokat, ellenőrzéseket végez, API-hívásokat indít, majd a kapott válasz alapján módosítja a felületet.

A választott technológiai alap előnye, hogy a rendszer fokozatosan bővíthető marad. Egy új funkció bevezetésekor nem feltétlenül szükséges a teljes alkalmazás szerkezetét átalakítani, hanem sok esetben egy meglévő nézet célzott bővítése vagy egy új oldal létrehozása is elegendő. Ez jól illeszkedik a WaccApp fejlesztési folyamatához, mert a rendszer nem egyszerre, egyetlen nagy felületként alakult ki, hanem egymásra épülő kommunikációs, kérdőíves, média- és szűrési funkciókból állt össze. A többnézetes szerkezet ugyanakkor csak akkor működik jól, ha az egyes oldalak hasonló működési elveket követnek, ezért a fejlesztésben fontos szerepet kaptak az ismétlődő űrlapminták, gombviselkedések, státuszjelzések és validációs megoldások. [R9]

\section{REST API-alapú kliens–szerver kommunikáció}

A WaccApp működésének egyik alapvető technológiai eleme a frontend és a backend közötti REST API-alapú kommunikáció. A kliensoldal HTTP-kéréseken keresztül továbbítja a felhasználó által megadott adatokat a háttérrendszer felé, a backend pedig ezek feldolgozása után strukturált, jellemzően JSON-formátumú választ ad vissza. A frontend ebből a válaszból nem technikai részletet jelenít meg közvetlenül, hanem a felület számára értelmezhető állapotot képez: sikeres műveletet, hibát, frissített listát vagy további felhasználói lépést. [R10], [R11], [R12]

A REST API-ra épülő működés azért illeszkedik jól a WaccApp felépítéséhez, mert a rendszer többféle műveletet kezel, és ezek eltérő adatokat igényelnek. Egy szöveges üzenetnél a címzett és az üzenettartalom a legfontosabb. Sablonos üzenetnél már a sablonnév, a nyelvi beállítás és a paraméterek is szerepet kapnak. Fájlküldésnél a kiválasztott állomány és annak kezelhetősége kerül előtérbe, szűrésnél pedig a megadott feltételek alapján kell értelmezhető eredménylistát előállítani. A frontend technológiai feladata ezekben az esetekben az, hogy az adott művelethez megfelelő adatcsomagot állítson össze, majd a választ a megfelelő nézetben dolgozza fel.

A WaccApp frontendjében az API-kommunikáció általában közvetlenül egy felhasználói művelethez kapcsolódik. A felhasználó kitölt egy űrlapot, kiválaszt egy sablont, kérdőívet vagy fájlt, beállít egy szűrési feltételt vagy időpontot, majd elindítja a műveletet. A JavaScript ekkor összegyűjti a szükséges adatokat, elvégzi az alapvető kliensoldali ellenőrzéseket, majd elküldi a kérést a megfelelő végpontra. A válasz megérkezése után a felület frissül: megjelenhet státuszszöveg, módosulhat egy gomb állapota, frissülhet egy lista, vagy hibaüzenet jelezheti a javítandó problémát.

A REST API használatának egyik előnye, hogy a frontend és a backend felelőssége viszonylag jól elkülöníthető. A frontend a felhasználói interakciókért, az adatbevitel elsődleges ellenőrzéséért, a megjelenítésért és a nézetfrissítésért felel. A backend ezzel szemben az adatbázis-kezelést, az üzleti logikát, az üzenetküldés tényleges végrehajtását és a WhatsApp Business Platformmal való kapcsolatot biztosítja. Ez a szétválasztás a fejlesztést és a hibakeresést is átláthatóbbá teszi, mert könnyebben megállapítható, hogy egy hiba a bevitelből, a kliensoldali adatfeldolgozásból, a végpont működéséből vagy a külső platform válaszából ered. [R10], [R11]

\section{Kliensoldali adatkezelés és nézetfrissítés}

A WaccApp frontendjében a kliensoldali adatkezelés fő szerepe az, hogy a felhasználó által megadott adatok ne rendezetlenül kerüljenek továbbításra a backend felé. A JavaScript először összegyűjti, csoportosítja és alapvetően ellenőrzi az adatokat, majd ezekből a művelet típusának megfelelő kérést állít össze. Ez a folyamat megjelenik az üzenetküldésnél, a sablonparaméterek kezelésénél, a kérdőívszerkesztésnél, a fájlkiválasztásnál, a szűrésnél és az időzített küldéseknél is.

Az \path{index.html} jó példa erre a működésre, mert ugyanazon indítófelületen többféle küldési mód találkozik. Egyszerű üzenetküldésnél a címzett és a szöveg a lényeges adat. Sablonos küldésnél a sablon neve és paraméterei is szükségesek. Kérdőívindításnál a kiválasztott kérdőív válik meghatározóvá, médiafájl esetén pedig a fájl és a hozzá kapcsolódó küldési adatok kerülnek előtérbe. A frontendnek ezeket az eltérő adatokat úgy kell kezelnie, hogy a felhasználó ne teljesen különálló rendszereket érzékeljen, hanem ugyanannak a kommunikációs felületnek a különböző működési módjait.

A kérdőívszerkesztésnél a kliensoldali adatkezelés összetettebbé válik. A \path{questionnaire.html} oldalon a kérdések, válaszlehetőségek és elágazások nem egymástól független mezők. Ezek kapcsolódó elemek, ahol egy válasz meghatározhatja a következő kérdést. A \path{template.html} és a kapcsolódó sablonkezelő nézetek esetében ehhez hozzájön a sablonos üzenetek szerkezete, a gombos válaszlogika és a következő lépések kezelése. Technológiai szempontból ez azt jelenti, hogy a frontend nemcsak adatokat olvas ki űrlapmezőkből, hanem egy kommunikációs folyamat szerkeszthető adatszerkezetét állítja elő.

A nézetfrissítés ehhez az adatkezeléshez kapcsolódik. Amikor egy API-hívás sikeresen lefut, a frontendnek módosítania kell az adott oldal tartalmát vagy állapotát. Ez lehet egy státuszüzenet megjelenítése, egy találati lista frissítése, egy beszélgetésfolyam újrarajzolása vagy egy mentési művelet eredményének jelzése. A \path{filter.html} például a megadott feltételek alapján lekért eredményeket jeleníti meg, míg a \path{conv.html} a kommunikációs előzményeket időrendben, beszélgetésszerű formában mutatja. A WaccApp frontendje ezért nem statikus dokumentumként működik, hanem olyan felületként, amely a felhasználói műveletek és az API-válaszok alapján folyamatosan változik.

\section{Aszinkron működés és kliensoldali állapotváltások}

A WaccApp több fontos művelete aszinkron módon működik, vagyis a felhasználói művelet elindítása és a végleges válasz megérkezése között idő telhet el. Ez technológiai szempontból azt jelenti, hogy a kliensoldalnak kezelnie kell azokat az átmeneti helyzeteket is, amikor a kérés már elindult, de az eredmény még nem ismert. Ilyen helyzet fordulhat elő üzenetküldésnél, sablonok betöltésénél, kérdőívek mentésénél, médiafájlok küldésénél, beszélgetések lekérésénél vagy szűrési eredmények megjelenítésénél. [R10], [R13]

A frontend működésében több alapállapot különíthető el. Nyugalmi állapotban a felület készen áll a használatra, a mezők kitölthetők, a gombok működnek, és a felhasználó elindíthatja a kívánt műveletet. Feldolgozás alatti állapot akkor jön létre, amikor a frontend már elküldte a kérést a backend felé, de a válasz még nem érkezett meg. Sikeres állapotban a rendszer visszajelzi, hogy a művelet rendben lezárult, hibás állapotban pedig a felületnek jeleznie kell, hogy a folyamat nem a várt eredménnyel fejeződött be.

Ebben a fejezetben az állapotkezelés technológiai oldala a lényeges. A kérdés nem pusztán az, hogy milyen szöveg jelenik meg a felhasználónak, hanem az, hogy a JavaScript mikor és hogyan módosítja a nézetet. Egy API-hívás elindításakor megváltozhat egy gomb állapota, megjelenhet egy feldolgozást jelző szöveg, vagy ideiglenesen módosulhat az adott nézet tartalma. A válasz megérkezése után a frontend új állapotba vált: frissíti a listát, kiírja a művelet eredményét, vagy hiba esetén javítható problémát jelez.

A több HTML-nézetből álló felépítés miatt különösen fontos volt, hogy az állapotváltások ne minden oldalon teljesen más logikát kövessenek. Az \path{index.html} küldési folyamatai, a \path{filter.html} lekérdezései, a \path{questionnaire.html} és a \path{template.html} mentési műveletei, valamint a \path{conv.html} adatbetöltései eltérő feladatokat látnak el, de technológiai szinten hasonló mintára épülnek: a felhasználó műveletet indít, a JavaScript API-kérést küld, majd a válasz alapján frissíti a felületet. Ez az egységes alapminta segít abban, hogy a rendszer többnézetes szerkezete ellenére technológiai szempontból követhető maradjon. [R13], [R14]

\section{Kliensoldali validációs és biztonsági alapok}

A WaccApp frontendjében a kliensoldali validáció célja az, hogy a nyilvánvalóan hibás vagy hiányos adatok lehetőség szerint már a kérés elküldése előtt felismerhetők legyenek. A tényleges ellenőrzés, jogosultságkezelés, üzleti logika és adatvédelmi védelem természetesen a backend feladata, de a frontend is hozzájárul ahhoz, hogy kevesebb hibás kérés jusson el a szerverhez. [R14]

Ez főként az űrlapok kezelésénél jelenik meg. Az \path{index.html} oldalon ellenőrizni kell, hogy a felhasználó megadta-e a címzettet, nem üres-e az üzenetszöveg, sablonküldés esetén rendelkezésre állnak-e a szükséges paraméterek, médiafájl küldésekor pedig kiválasztott-e megfelelő fájlt. Ezek az ellenőrzések egyszerűnek tűnhetnek, de a működés szempontjából fontosak, mert a felhasználó gyorsabban javíthat, és a backend nem kap olyan kéréseket, amelyek már kliensoldalon is egyértelműen hibásnak tekinthetők.

A kérdőíves moduloknál a validáció összetettebb szerepet kap. A \path{questionnaire.html} esetében nemcsak azt kell figyelni, hogy a kérdések és válaszok ki vannak-e töltve, hanem azt is, hogy a válaszokhoz tartozó következő lépések értelmezhetők-e. A \path{template.html} oldalon hasonló probléma jelenik meg a sablonos kérdőívfolyamatoknál, ahol a gombos válaszokhoz tartozó folytatásokat következetesen kell megadni. Ha egy kérdőívláncban hiányzik egy következő elem, az később nem egyszerű megjelenítési hiba lesz, hanem az egész kommunikációs folyamat megszakadásához is vezethet.

Biztonsági szempontból a frontend szerepe korlátozottabb, de nem elhanyagolható. A WaccApp kliensoldala meghatározott backend végpontokkal kommunikál, és a felhasználói műveleteket ezeken keresztül továbbítja. A CORS-szabályozás és a végleges jogosultsági ellenőrzés szerveroldali feladat, de frontendoldalon is fontos, hogy a rendszer előre meghatározott működési környezetre épüljön, és a felület ne sugalljon olyan műveleti szabadságot, amelyet a backend vagy a WhatsApp Business Platform szabályai később nem engednek meg. A kliensoldali validáció tehát technológiai előszűrésként jelenik meg. A hibás inputok kiszűrése, a kötelező mezők kezelése, a fájltípusok ellenőrzése és a kérdőívlogika alapvető vizsgálata mind azt szolgálja, hogy a frontend rendezettebb, előkészítettebb kéréseket küldjön tovább a backend felé. [R11]

\section{A választott technológiai megközelítés előnyei és kompromisszumai}

A WaccApp frontendje több különálló HTML-nézetből és hozzájuk kapcsolódó kliensoldali működési logikából épül fel. Ennek a megközelítésnek az egyik legnagyobb előnye az átláthatóság. Az egyes funkciók elkülöníthető oldalakra kerültek, így a fejlesztés és a hibakeresés könnyebben követhető maradt. Ha a szűrési logikában jelentkezik probléma, akkor elsősorban a \path{filter.html} és a hozzá kapcsolódó JavaScript vizsgálata kerül előtérbe. Ha a beszélgetésmegjelenítés hibás, akkor a \path{conv.html} és az ahhoz kapcsolódó adatlekérési folyamat válik elsődlegessé.

A többoldalas, kliensoldali JavaScriptre épülő megoldás a projekt méretéhez is illeszkedett. Egy összetettebb frontendkeretrendszer bizonyos szempontból kényelmesebb lehetett volna, például központi állapotkezelés vagy újrahasznosítható komponensek esetén. A WaccApp jelenlegi funkcióihoz azonban nem volt feltétlenül indokolt egy nagyobb frontend-infrastruktúra bevezetése. A cél nem az volt, hogy a rendszer technológiailag minél bonyolultabb legyen, hanem az, hogy a választott megoldás érthetően támogassa az üzenetküldést, a kérdőívkezelést, a médiafájlok kezelését, az időzítést és a visszakeresést. [R11], [R14]

A fejlesztés korai szakaszában, illetve a témaleírás technológiai irányainak megfogalmazásakor felmerült egy keretrendszer-alapú, komponensvezérelt frontend lehetősége is. Ilyen megközelítésben a React és egy utility-first szemléletű stíluskezelés, például Tailwind CSS is alkalmas lett volna egy egységes komponensrendszer kialakítására. A WaccApp tényleges fejlesztése során azonban a funkciók természete inkább több, jól elkülöníthető HTML-nézetet indokolt. Az üzenetküldés, a beszélgetésnézet, a szűrés, a kérdőívszerkesztés és a sablonkezelés önálló felületi helyzeteket képviseltek, ezért a többoldalas HTML/CSS/JavaScript-alapú megoldás átláthatóbbnak és gyorsabban módosíthatónak bizonyult. Ez a döntés ugyan kevesebb központi komponenslogikát eredményezett, de a projekt méretéhez, fejlesztési folyamatához és frontendfókuszához jobban illeszkedett.

A választott megoldásnak ugyanakkor vannak kompromisszumai. Egy modern egyoldalas alkalmazás vagy komponensalapú frontendkeretrendszer előnyösebb lehetne akkor, ha a WaccApp sok közös globális állapotot, valós idejű nézetváltást, összetettebb komponenshierarchiát vagy nagyobb mértékű újrafelhasználhatóságot igényelne. Ilyen esetben egy központibb frontendarchitektúra jobban kezelné a megosztott állapotokat és az ismétlődő felületi elemeket.

A WaccApp jelenlegi fejlesztési szintjén azonban a többnézetes felépítés megfelelő kompromisszumnak bizonyult. Egyszerűen követhető, jól dokumentálható, és a fő funkciók külön-külön is értelmezhetők maradnak. Az egységes működést ebben az esetben nem egy keretrendszer biztosítja, hanem a fejlesztés során következetesen alkalmazott vizuális minták, validációs megoldások és állapotkezelési elvek. Ez a megközelítés illeszkedik a projekt gyakorlati céljához: egy olyan frontend kialakításához, amely a WhatsApp Business API-ra épülő kommunikációs folyamatokat átlátható webes felületen teszi kezelhetővé.

\section{Futtatási és telepítési szempontok}

A WaccApp technológiai hátterének gyakorlati oldalához a futtatási és telepítési környezet egyszerűsége is hozzátartozik. Mivel a frontend különálló HTML-oldalakból és hozzájuk kapcsolódó kliensoldali állományokból épül fel, a rendszer kiszolgálása jól illeszthető a backend működéséhez. A frontend elhelyezhető a backend által kiszolgált nyilvános mappában, így a felhasználó egy webes környezeten belül éri el a kommunikációs funkciókat.

Fejlesztői szempontból ez a megoldás azért is kedvező, mert az egyes módosítások viszonylag gyorsan érvényesíthetők, az oldalak önállóan tesztelhetők, és a hibák könnyebben elkülöníthetők. Egy szűrési hiba vizsgálatakor például nem kell azonnal a teljes frontend működését áttekinteni, hanem elsőként a szűrőnézet és az ahhoz kapcsolódó adatlekérés ellenőrizhető. Ugyanez igaz a beszélgetésmegjelenítésre, a kérdőívszerkesztésre vagy a küldési folyamatokra is.

A futtatási szempontoknál ugyanakkor fontos kiemelni, hogy a WaccApp frontendje nem egymástól független weboldalak gyűjteménye. Az oldalak a szerveroldali végpontokkal együtt alkotnak teljes rendszert. A felhasználó ebből elsősorban azt érzékeli, hogy elérhető nézeteket, kitölthető űrlapokat, visszajelzéseket, szűrési eredményeket és követhető kommunikációs folyamatokat kap. A technológiai háttér ebben az értelemben akkor tölti be a szerepét, ha a HTML-, CSS- és JavaScript-alapú frontend, valamint a backend API-kapcsolatai együtt stabil és használható felületet eredményeznek.
